% from ega4-16-fr.tex
\oldpage[IV-4]{30}
d'où l'on déduit, en vertu de (16.5.10.1), un homomorphisme canonique
(Bourbaki, \emph{Alg.}, chap. II, 3\textsuperscript{e} éd., \S~5,
n\textsuperscript{o} 3)
\[
  \mathfrak{G}_{X/S}\otimes_{\mathcal{O}_S}\mathcal{O}_{S'}
  \longrightarrow\mathfrak{G}_{X'/S'}
  \tag{16.5.10.2}
  \label{IV.16.5.10.2}
\]
qui en général n'est ni injectif ni surjectif. Toutefois:
\end{env}

\begin{proposition}[16.5.11]
\label{IV.16.5.11}
\begin{enumerate}
\item[(i)] Si $g:S'\to S$ est un morphisme plat et si $f$ est localement de
type fini (resp.\ localement de présentation finie), l'homomorphisme
(16.5.10.2) est injectif (resp.\ bijectif).
\item[(ii)] Si $\Omega_{X/S}^1$ est un $\mathcal{O}_X$-Module localement libre
de type fini, l'homomorphisme (16.5.10.2) est bijectif.
\end{enumerate}
\end{proposition}

En effet, l'assertion (ii) résulte de Bourbaki, \emph{Alg.}, chap. II,
3\textsuperscript{e} éd., \S~5, n\textsuperscript{o} 3, prop. 7. L'assertion
(i) résulte de même de Bourbaki, \emph{Alg. comm.}, chap. I, \S~2,
n\textsuperscript{o} 10, prop. 11 et du fait que si $f$ est localement de type
fini (resp.\ localement de présentation finie), $\Omega_{X/S}^1$ est un
$\mathcal{O}_X$-Module de type fini (resp.\ de présentation finie) ((16.3.9)
et (16.4.22)).

\begin{env}[16.5.12]
\EGAReaderTarget{ega4.c04.statement.16-5-12}{IV.16.5.12}
Puisque $\Omega_{X/S}^1$ est un $\mathcal{O}_X$-Module quasi-cohérent, on peut
considérer le fibré vectoriel sur $X$ défini par $\Omega_{X/S}^1$
(II, 1.7.8)
\[
  T_{X/S}=\mathbf{V}(\Omega_{X/S}^1)
  \tag{16.5.12.1}
  \label{IV.16.5.12.1}
\]
qu'on appelle \emph{fibré tangent de $X$ relativement à $S$}. On a donc une
bijection canonique (II, 1.7.9)
\[
  \Gamma(T_{X/S}/S)\xrightarrow{\sim}
  \operatorname{Hom}_{\mathcal{O}_X}(\Omega_{X/S}^1,\mathcal{O}_X)
  =\Gamma(X,\mathfrak{G}_{X/S})
\]
par définition de $\mathfrak{G}_{X/S}$, et on peut dans cet isomorphisme
remplacer $X$ par un ouvert quelconque $U$ de $X$; on peut donc dire que le
\emph{faisceau tangent} de $X$ relativement à $S$ est isomorphe au
\emph{faisceau des germes de $S$-sections} du fibré tangent de $X$
relativement à $S$. Si $f:X\to Y$ est un $S$-morphisme, on a vu (16.4.19)
qu'on a un homomorphisme canonique
$f_{X/Y/S}:f^*(\Omega_{Y/S}^1)\to\Omega_{X/S}^1$, ce qui donne, compte tenu de
ce que
\[
  \mathbf{V}(f^*(\Omega_{Y/S}^1))
  =\mathbf{V}(\Omega_{Y/S}^1)\times_Y X\qquad\text{(II, 1.7.11),}
\]
un $X$-morphisme
$T_{X/S}(f):T_{X/S}\to T_{Y/S}\times_Y X$. Si $g:Y\to Z$ est un second
$S$-morphisme, on a
$T_{X/S}(g\circ f)=(T_{Y/S}(g)\times1_X)\circ T_{X/S}(f)$ (0, 20.5.4.1).

Il résulte de (16.5.10.1) et de (II, 1.7.11) que pour tout changement de base
$g:S'\to S$, on a un isomorphisme canonique
\[
  T_{X'/S'}\xrightarrow{\sim}T_{X/S}\times_S S'
  =T_{X/S}\times_X X'.
  \tag{16.5.12.2}
  \label{IV.16.5.12.2}
\]
\end{env}

\begin{env}[16.5.13]
\EGAReaderTarget{ega4.c04.statement.16-5-13}{IV.16.5.13}
Pour tout point $x\in X$, on appelle \emph{espace tangent à $X$ au point $x$}
(relativement à $S$) l'ensemble des points de la fibre
$T_{X/S}\times_X\operatorname{Spec}(k(x))$ \emph{rationnels sur $k(x)$}, donc
l'ensemble
\[
  T_{X/S}(x)=
  \operatorname{Hom}_{k(x)}
  (\Omega_{X/S}^1\otimes_{\mathcal{O}_x}k(x),k(x))
  \tag{16.5.13.1}
  \label{IV.16.5.13.1}
\]
qui est le \emph{dual} du $k(x)$-espace vectoriel
$\Omega_{\mathcal{O}_x/\mathcal{O}_s}^1/
\mathfrak{m}_x\cdot\Omega_{\mathcal{O}_x/\mathcal{O}_s}^1$. Lorsque
$\Omega_{X/S}^1$ est un $\mathcal{O}_X$-Module de \emph{type fini},
$T_{X/S}(x)$ est donc un espace vectoriel de rang fini sur $k(x)$, et pour
tout
